\begin{center}
  {\zihao{2}\heiti 嵌入式社区养老服务站的建设与优化问题}
\end{center}

\section*{摘要}

\noindent
嵌入式社区养老服务站建设需要同时回应老人规模增长、服务需求差异、站点建设预算、服务半径、服务能力、价格补贴和运营风险等约束。针对 B 题要求，本文围绕“需求预测--消费约束修正--选址规模优化--补贴导向定价--可及性评价--灵敏度分析”的主线，建立三状态老人数量递推模型、收费服务等比例削减模型、站点--规模组合枚举与预算剪枝模型、满意度--利用率--容量--分配迭代算法、补贴导向统一定价模型和三类老人服务可及性评价模型。模型中将紧急救助作为公益免费服务处理，不参与消费约束、不计政府补贴，但计入直接成本和服务能力占用。

\noindent
针对问题一，本文预测第 5 年末老人总数为 $7578.41$ 人，其中自理、半失能和失能老人分别为 $4981.93$ 人、$1470.85$ 人和 $1125.63$ 人。第 5 年末理论月服务需求为 $262609.62$ 次/月，经消费约束修正后实际月服务需求为 $247753.32$ 次/月；紧急救助需求保持不削减，收费服务按同一比例修正，从而在满足消费预算的同时保持服务结构相对稳定。针对问题二，本文在 $120$ 万元建设预算和 $1000$ 米服务半径下枚举 $4^{10}$ 个站点--规模组合，得到最优方案为建设 $C$ 小型站、$D$ 大型站和 $G$ 大型站，总建设成本 $108$ 万元。该方案题目口径服务覆盖率和地理覆盖率均为 $1.0000$，有效服务覆盖率为 $0.8530$，需求满足率为 $0.8531$，老人加权平均满意度为 $0.8622$；年度利润仅作为基准价格下的测算结果，不参与选址排序。

\noindent
针对问题三，本文在固定问题二站点方案的基础上进行补贴导向定价搜索。最终统一服务价格为助餐 $7$ 元、日间照料 $18$ 元、上门护理 $29$ 元、康复理疗 $27$ 元、助浴 $21$ 元，紧急救助保持 $0$ 元。定价后老人年度支付额为 $35647892.66$ 元，年度政府补贴为 $2263000.00$ 元；三处站点题目口径利润率区间为 $[-12.20\%,7.72\%]$，满足不超过 $8\%$ 的利润率上限。由于 $C$ 小型站年度净收益为 $-98068.66$ 元，说明方案存在局部保本风险，需要财政统筹或运营效率改进。三类老人经济可及性均为 $1.0000$，综合可及性约为 $0.929$，表明补贴定价后类型间服务获得能力差异较小，但服务满足维度仍是进一步优化重点。

\noindent
针对问题四，本文对老人增长率、健康状态转移概率、固定管理成本和建设预算变化进行完整重求解。结果表明，老人增长率升至 $8\%$ 时站点位置发生调整但综合变化度仅为 $0.2133$；转移概率变化和固定成本上升分别使年度净收益降至 $-148.90$ 万元和 $-183.59$ 万元，其中固定成本扰动的综合变化度最高，为 $2.0227$。当建设预算提高至 $140$ 万元时，方案扩展为 $A$ 大型、$E$ 中型、$G$ 大型和 $I$ 小型四站，有效覆盖率提高至 $0.9017$、满意度提高至 $0.9769$，但年度净收益仍为 $-70.54$ 万元。由此可见，预算增加能够改善覆盖和满意度，但不能直接等同于运营更优，仍需结合补贴压力和保本风险筛选实施方案。

\noindent\textbf{关键词：}社区养老服务站；服务需求预测；选址优化；补贴定价；可及性评价
